% ============================================================================
% WORKBOOK — Section 5.4: Solving Multi-Step Problems with Rational Numbers
% CCSS 7.EE.B.3
% Practice-focused reinforcement for the study guide topic
% ============================================================================
\section{Solving Multi-Step Problems with Rational Numbers}
\topicTitle{Solving Multi-Step Problems with Rational Numbers}

\begin{quickReview}{Equations with Fractions and Decimals}
When equations have fractions or decimals, two helpful strategies:
\begin{itemize}[leftmargin=*]
    \item \textbf{Clear fractions:} Multiply every term by the LCD (least common denominator).
    
    $\frac{x}{3} + \frac{1}{2} = 5$ → multiply everything by $6$: $2x + 3 = 30$.
    \item \textbf{Work with decimals directly:} Use standard arithmetic.
    
    $0.25x + 1.5 = 4$ → subtract $1.5$: $0.25x = 2.5$ → $x = 10$.
\end{itemize}

\medskip
\textbf{Tip:} Before solving, \textbf{estimate} the answer. If $0.5x + 3 = 10$, then $x \approx 14$. If you get $140$, check your decimal work!
\end{quickReview}

% ── Warm-Up ──────────────────────────────────────────────────────────────────
\begin{practiceBox}[funGreen]{\faSun~Warm-Up}

\practiceHeader[funGreen]{Quick Decimal and Fraction Practice}

\begin{multicols}{2}
\prob $0.5 \times 12 = $ \answerBlank[2cm]
\answerExplain{$6$}{Multiply: $0.5 \times 12 = 6$. Half of $12$ is $6$.}

\prob $\dfrac{3}{4} \times 20 = $ \answerBlank[2cm]
\answerExplain{$15$}{Multiply: $\frac{3}{4} \times 20 = \frac{60}{4} = 15$.}

\prob $2.4 \div 0.6 = $ \answerBlank[2cm]
\answerExplain{$4$}{Divide: $2.4 \div 0.6 = 4$. You can also think: $0.6 \times 4 = 2.4$.}

\prob $\dfrac{2}{3} + \dfrac{1}{6} = $ \answerBlank[2cm]
\answerExplain{$\frac{5}{6}$}{Find a common denominator: $\frac{4}{6} + \frac{1}{6} = \frac{5}{6}$.}

\prob $7.5 - 3.8 = $ \answerBlank[2cm]
\answerExplain{$3.7$}{Subtract: $7.5 - 3.8 = 3.7$.}

\prob LCD of $3$ and $4$ is \answerBlank[2cm]
\answerExplain{$12$}{The least common denominator of $3$ and $4$ is $12$, since $12$ is the smallest number divisible by both.}
\end{multicols}
\end{practiceBox}

% ── Practice A: Decimal Equations ────────────────────────────────────────────
\begin{practiceBox}{Solving Equations with Decimals}

\practiceHeader{Solve each equation.}

\prob $0.4x + 3 = 7$ \answerBlank[3cm]
\answerExplain{$x = 10$}{Subtract $3$: $0.4x = 4$. Divide by $0.4$: $x = 10$. Check: $0.4(10) + 3 = 4 + 3 = 7$ \checkmark.}

\prob $1.5n - 4.5 = 6$ \answerBlank[3cm]
\answerExplain{$n = 7$}{Add $4.5$: $1.5n = 10.5$. Divide by $1.5$: $n = 7$. Check: $1.5(7) - 4.5 = 10.5 - 4.5 = 6$ \checkmark.}

\prob $-0.2y + 5 = 3.4$ \answerBlank[3cm]
\answerExplain{$y = 8$}{Subtract $5$: $-0.2y = -1.6$. Divide by $-0.2$: $y = 8$. Check: $-0.2(8) + 5 = -1.6 + 5 = 3.4$ \checkmark.}

\prob $0.25a + 0.75 = 2$ \answerBlank[3cm]
\answerExplain{$a = 5$}{Subtract $0.75$: $0.25a = 1.25$. Divide by $0.25$: $a = 5$. Check: $0.25(5) + 0.75 = 1.25 + 0.75 = 2$ \checkmark.}

\prob $3.2x - 1.6 = 8$ \answerBlank[3cm]
\answerExplain{$x = 3$}{Add $1.6$: $3.2x = 9.6$. Divide by $3.2$: $x = 3$. Check: $3.2(3) - 1.6 = 9.6 - 1.6 = 8$ \checkmark.}

\end{practiceBox}

% ── Practice B: Fraction Equations ───────────────────────────────────────────
\begin{practiceBox}[funTeal]{Solving Equations with Fractions}

\practiceHeader[funTeal]{Clear the fractions by multiplying by the LCD, then solve.}

\prob $\dfrac{x}{4} + 3 = 10$ \answerBlank[3cm]
\answerExplain{$x = 28$}{Subtract $3$: $\frac{x}{4} = 7$. Multiply by $4$: $x = 28$. Check: $\frac{28}{4} + 3 = 7 + 3 = 10$ \checkmark.}

\prob $\dfrac{2n}{3} - 5 = 1$ \answerBlank[3cm]
\answerExplain{$n = 9$}{Add $5$: $\frac{2n}{3} = 6$. Multiply by $3$: $2n = 18$. Divide by $2$: $n = 9$. Check: $\frac{2(9)}{3} - 5 = 6 - 5 = 1$ \checkmark.}

\prob $\dfrac{x}{2} + \dfrac{x}{4} = 9$ \answerBlank[3cm]
\answerExplain{$x = 12$}{Multiply every term by $4$ (LCD): $2x + x = 36$. Combine: $3x = 36$. Divide by $3$: $x = 12$. Check: $\frac{12}{2} + \frac{12}{4} = 6 + 3 = 9$ \checkmark.}

\prob $\dfrac{3y}{5} + 2 = 8$ \answerBlank[3cm]
\answerExplain{$y = 10$}{Subtract $2$: $\frac{3y}{5} = 6$. Multiply by $5$: $3y = 30$. Divide by $3$: $y = 10$. Check: $\frac{3(10)}{5} + 2 = 6 + 2 = 8$ \checkmark.}

\prob $\dfrac{m}{6} - \dfrac{1}{3} = \dfrac{3}{2}$ \answerBlank[3cm]
\answerExplain{$m = 11$}{Multiply every term by $6$ (LCD): $m - 2 = 9$. Add $2$: $m = 11$. Check: $\frac{11}{6} - \frac{1}{3} = \frac{11}{6} - \frac{2}{6} = \frac{9}{6} = \frac{3}{2}$ \checkmark.}

\end{practiceBox}

% ── Practice C: True or False ────────────────────────────────────────────────
\begin{practiceBox}[funPurple]{True or False?}

\trueOrFalse{To solve $\frac{x}{3} + \frac{1}{6} = 2$, you can multiply every term by $6$ to clear the fractions.}
\answerExplain{True}{The LCD of $3$ and $6$ is $6$. Multiplying every term by $6$ gives $2x + 1 = 12$, which has no fractions and is easier to solve.}

\trueOrFalse{The solution to $0.5x + 2 = 7$ is $x = 18$.}
\answerExplain{False}{Subtract $2$: $0.5x = 5$. Divide by $0.5$: $x = 10$, not $18$. The student likely made an arithmetic error.}

\trueOrFalse{You should always clear fractions before solving. Working with fractions directly is not a valid method.}
\answerExplain{False}{Both methods are valid. Sometimes working with fractions directly is just as fast. Clearing fractions is a helpful strategy but not the only one.}

\end{practiceBox}

% ── Practice D: Word Problems ────────────────────────────────────────────────
\begin{practiceBox}[funBlue]{\faGlobe~Word Problems}

\wordProblem{A recipe uses $\frac{3}{4}$ cup of sugar per batch plus $\frac{1}{2}$ cup of vanilla sugar (added once). You used $3\frac{1}{2}$ cups of sugar total. How many batches did you make?}{batches}
\answerExplain{$4$ batches}{Let $b =$ batches. Equation: $\frac{3}{4}b + \frac{1}{2} = 3\frac{1}{2}$. Subtract $\frac{1}{2}$: $\frac{3}{4}b = 3$. Multiply by $\frac{4}{3}$: $b = 4$. Check: $\frac{3}{4}(4) + \frac{1}{2} = 3 + \frac{1}{2} = 3\frac{1}{2}$ \checkmark.}

\wordProblem{A store discounts jeans by $\frac{1}{4}$ of the original price, then takes off another $\$5$. The final price is $\$25$. What was the original price?}{dollars}
\answerExplain{$\$40$}{Let $p$ = original price. After discount: $p - \frac{p}{4} = \frac{3p}{4}$. Then subtract $\$5$: $\frac{3p}{4} - 5 = 25$. Add $5$: $\frac{3p}{4} = 30$. Multiply by $\frac{4}{3}$: $p = 40$.}

\wordProblem{Sam earned $\$45$ for working $3.5$ hours plus a $\$10$ bonus. What is his hourly rate?}{dollars per hour}
\answerExplain{$\$10$ per hour}{Let $r =$ hourly rate. Equation: $3.5r + 10 = 45$. Subtract $10$: $3.5r = 35$. Divide by $3.5$: $r = 10$. Check: $3.5(10) + 10 = 35 + 10 = 45$ \checkmark.}

\wordProblem{A tank holds $\frac{2}{3}$ of its capacity. After adding $15$ gallons, it is $\frac{5}{6}$ full. What is the tank's capacity?}{gallons}
\answerExplain{$90$ gallons}{Let $c =$ capacity. Equation: $\frac{2}{3}c + 15 = \frac{5}{6}c$. Subtract $\frac{2}{3}c$: $15 = \frac{5}{6}c - \frac{4}{6}c = \frac{1}{6}c$. Multiply by $6$: $c = 90$ gallons.}

\end{practiceBox}

% ── Challenge ────────────────────────────────────────────────────────────────
\begin{challengeBox}

\prob Solve $\dfrac{2x - 3}{5} + 1 = 4$.
\answerExplain{$x = 9$}{Subtract $1$: $\frac{2x - 3}{5} = 3$. Multiply by $5$: $2x - 3 = 15$. Add $3$: $2x = 18$. Divide by $2$: $x = 9$. Check: $\frac{2(9) - 3}{5} + 1 = \frac{15}{5} + 1 = 3 + 1 = 4$ \checkmark.}

\prob A bookstore sells hardcovers for $\$12.50$ each and paperbacks for $\$7.50$ each. You buy $2$ hardcovers and some paperbacks. Your total before tax is $\$47.50$. How many paperbacks did you buy?
\answerExplain{$3$ paperbacks}{Hardcovers: $2 \times 12.50 = \$25$. Equation: $25 + 7.50p = 47.50$. Subtract $25$: $7.50p = 22.50$. Divide by $7.50$: $p = 3$. Check: $25 + 3(7.50) = 25 + 22.50 = 47.50$ \checkmark.}

\end{challengeBox}

\encouragement{Fantastic work with fractions and decimals in equations! You are building real problem-solving power!}
